\chapter{Comparing methods by what they estimate}
\label{sw:interfaces}

\section{The decision determines the comparison}

A bank does not need an abstract winner among causal libraries. It needs to
know whether a particular computation answers its particular question. A
funding desk might ask how changing a quoted rate affects twelve-month retained
balances. A servicing team might ask when to contact an account whose financial
condition changes after earlier contacts. A model-risk team might ask whether
two plausible causal diagrams imply different intervention decisions. These
questions require different information, estimands and outputs, even when the
software accepts the same rectangular dataset.

The useful comparison therefore starts with a scientific interface. Let
\(q=(\mathcal P,A,Y,h,s)\) denote a target population \(\mathcal P\), an intervention
\(A\), an outcome \(Y\), a horizon \(h\), and a scale \(s\). A request also supplies
an information history, a causal model or model class, a sampling design and
explicit assumptions. An implementation may identify a functional of the
observed distribution, estimate that functional, characterize uncertainty, or
report that its method does not support the request. These are different
operations. A scalar with a standard error is inadequate if the scientifically
appropriate answer is an equivalence class of graphs, a disconnected set of
effects, or a symbolic expression for which no numerical estimator has been
provided.

This distinction organizes the selected software comparisons developed here.
The evidence is package-authored documentation and selected implementation
paths at explicitly recorded revisions. A documented option is compared with
the relevant source where available. These inspections establish what was
read and what interfaces expose; they do not establish installation success,
numerical parity, statistical calibration or suitability for confidential bank
data. The banking calculations are project examples with specified mechanisms,
not measurements from those packages or from a bank.

The language of a package matters because it constrains integration, but it is
not a scientific ranking. R offers several specialized design and longitudinal
interfaces. Python offers graph, discovery and orchestration interfaces.
Tetrad supplies a Java implementation family. aGrUM exposes C++ causal-model
and identification interfaces with a related Python wrapper. A wrapper can
reduce deployment effort while changing exception behavior or result types;
a port can make deployment more uniform while introducing a second numerical
implementation to validate. These tradeoffs are assessed at the operation and
result boundary, rather than by counting languages or API names.

\section{A graph can be known while an effect remains unresolved}

Consider a deliberately simple project example involving a standardized quoted
rate \(X\) and a standardized balance response \(Y\). Two linear Gaussian
structural models are
\begin{align}
 \mathcal M_1 &: X=U_X,\qquad Y=0.4X+U_Y,
    &\operatorname{var}(U_X)&=1,&\operatorname{var}(U_Y)&=0.84,\\
 \mathcal M_2 &: Y=V_Y,\qquad X=0.4Y+V_X,
    &\operatorname{var}(V_Y)&=1,&\operatorname{var}(V_X)&=0.84.
\end{align}
All disturbances are centered and independent within their model. Both models
generate a bivariate normal distribution with unit marginal variances and
covariance \(0.4\). Under \(\mathcal M_1\), increasing \(X\) by one standardized
unit through an intervention changes the mean of \(Y\) by \(0.4\). Under
\(\mathcal M_2\), replacing the equation for \(X\) leaves \(Y\) unchanged. The
observational covariance does not choose between these answers. The example
is an algebraic construction, not an assertion that either model describes
actual funding behavior.

A graph interface should therefore preserve the distinction between an
asserted causal DAG and a class learned from observations. The inspected
Tetrad PC documentation states causal sufficiency, acyclicity and faithfulness,
uses an externally supplied conditional-independence test, and describes CPDAG
output. Its FCI search path contains circle endpoints, possible-d-separation
processing and PAG-oriented logic. An inconsistent class-level description
calling that output a CPDAG cannot serve as the interchange specification.
The algorithm path and the actual endpoint representation must be checked
together.\cite{v12_sw_tetrad}

The inspected causal-learn PC interface makes the dependence on procedural
choices explicit: it accepts a conditional-independence test, a test level,
a stabilized skeleton option, collider-orientation rules, conflict priorities
and background knowledge. Its graph matrix can encode two arrowheads, and
one documented conflict priority deliberately permits that outcome. Thus a
bidirected mark produced during a PC conflict must not automatically be
translated into an FCI-style statement about latent confounding. The FCI
interface separately documents tails, arrowheads and circles and marks its
result as a PAG. Both endpoint values and the method that produced them
belong in the result.\cite{v12_sw_causallearn}

Search bounds also belong in that result. causal-learn exposes a conditioning
bound for PC and adjacency-depth and discriminating-path bounds for FCI. These
can change the computation being performed. A time limit or restricted search
is not merely a performance annotation if it changes which separating sets or
orientation rules can be examined. The proposed adapter must retain these
settings and distinguish completed bounded search from interrupted search.
The conditional-independence test level is a setting of that test; no whole-graph
confidence level follows from the presence of a single parameter named alpha.

For a banking application, background knowledge should specify the proposition
being added and its provenance. An accounting identity, a verified temporal
ordering and a language model's directional suggestion carry different
evidentiary weight. A useful graph object retains alternatives until additional
information resolves them. Exporting the first available DAG would convert
computational convenience into an unrecorded identifying assumption.

\hypertarget{java-discovery-the-statistical-contract-before-the-engine}{%
\section{Java discovery: the statistical contract before the
engine}\label{java-discovery-the-statistical-contract-before-the-engine}}

\label{exec13:java-discovery}

For a funding desk, an attractive discovery interface accepts a table of
account histories and returns a diagram. The hard question is what that
diagram means. An undirected link can represent unresolved orientation,
a selected classifier can point an outcome toward its predictors, and a
directed optimization result can depend on its variable scale. A Java
boundary must preserve these differences before the proposal can use its
output in an intervention analysis. The comparison therefore follows the
statistical operation from input preparation through search, graph
representation and failure handling.

The primary comparison covers selected Tetrad sources at the 22
September 2026 revision, the discontinued Weka 3.8 mirror at its 10
August 2022 revision, Aitiotita at its 9 September 2023 revision, and a
Banjo derivative at its 13 August 2026 revision. These are different
kinds of evidence. The Weka mirror expressly directs users to the
official Waikato repository, while the Banjo derivative supplies an
application and selected inherited implementation files rather than an
authenticated original distribution. Revision-specific findings cannot
establish what every current release provides. The comparison below uses
inspected technical text; it contains no execution results from these
Java packages.
\citep{exec13_java_tetrad,exec13_java_weka,exec13_java_aitiotita,exec13_java_banjo}

\hypertarget{four-search-families-answer-different-questions}{%
\subsection{Four search families answer different
questions}\label{four-search-families-answer-different-questions}}

Constraint-based discovery searches for conditional independences and
then orients what the assumed graph family and those independences
permit. Tetrad's PC documentation specifies sufficiency, acyclicity and
faithfulness, accepts a separate independence test, and returns a CPDAG.
Sufficiency excludes unmeasured common causes in that model;
faithfulness rules out distributional cancellations that hide graph
connections. Neither property follows from a spreadsheet containing many
variables. A funding table that omits a relationship manager's private
assessment may therefore fall outside the declared PC model even when
every measured field is complete. \citep{exec13_java_tetrad}

Score-based search compares statistical representations using a selected
score. The inspected FGES constructor takes that score as an input and
the search alternates forward and backward passes, with processing for a
completed partially directed graph. The inspected SEM-BIC score is
documented for an IID linear Gaussian model, with \[
 S(G)=2\ell_G-c k_G\log n.
\] Here \(\ell_G\) is the fitted log likelihood, \(k_G\) its parameter
count, \(n\) the sample size and \(c\) a penalty multiplier, initialized
to one in the reviewed score. Increasing \(c\) changes the preference
for added parameters; it does not create a known probability that every
returned edge is correct. Accounts repeated across months do not become
IID observations because the score accepts the number of table rows as
\(n\). A proposal using that score must first choose the observation
unit, sampling regime and treatment of dependence.
\citep{exec13_java_tetrad}

Functional discovery adds assumptions about the structural equations and
their disturbances. In Tetrad's inspected DirectLingam path, variables
are standardized, an exogeneity criterion chooses a next variable,
remaining variables are residualized, and a supplied score enters a
subsequent grow-shrink parent selection. Thus the selected order and the
selected parent sets are distinct parts of the computation. A directed
result carries the consequences of that functional model. Counting its
directed edges against the partially oriented output of PC would reward
an additional assumption without asking whether the banking mechanism
supports it. \citep{exec13_java_tetrad}

Continuous optimization offers another implementation route. The
inspected Tetrad DAGMA constructor starts from the data covariance,
initializes a sparsity penalty of 0.05 and a weight threshold of 0.1,
and defaults to a CPDAG result. Its conversion routine thresholds
absolute weights, adds an edge from variable \(i\) to variable \(j\) for
a retained \(W_{ij}\), and optionally applies a DAG-to-CPDAG
transformation. The graph output consequently does not itself retain the
signed coefficient matrix. The optimization schedule and a valid acyclic
output are separate from an empirical guarantee that the selected causal
structure is correct. This is a source-level description of the Java
path, not a new theoretical review or numerical validation of DAGMA.
\citep{exec13_java_tetrad}

For the proposed library these families should remain explicit choices.
A request should name the observation law, graph class, allowed latent
structure, functional restrictions, test or score, search bounds, and
the role of prior knowledge. The returned object should name its class
as well. A selected DAG, a CPDAG and a PAG cannot share a field whose
only interpretation is ``the causal graph.'\,' The preceding
observational-equivalence example explains why this is consequential for
the effect; the Java implementations show where that distinction can be
lost in an adapter.

\hypertarget{conditional-independence-failure-is-part-of-the-scientific-result}{%
\section{Conditional-independence failure is part of the scientific
result}\label{conditional-independence-failure-is-part-of-the-scientific-result}}

\hypertarget{what-fishers-transform-requires}{%
\subsection{What Fisher's transform
requires}\label{what-fishers-transform-requires}}

In the inspected Tetrad Fisher-Z computation, the conditional sample
correlation \(r_{XY\mid Z}\) is transformed using \[
 z_F=\sqrt{n_{\mathrm{eff}}-3-|Z|}\,
       \operatorname{atanh}(|r_{XY\mid Z}|),\qquad
 p=2\{1-\Phi(z_F)\}.
\] The normal-tail calculation is the implementation's reference test,
with \(n_{\mathrm{eff}}\) the configured effective sample size and
\(|Z|\) the conditioning-set cardinality. It does not supply an
automatic solution for serially dependent customer histories, heavy
tails or selection into an observed table. A user who changes the
effective sample size has changed a calculation; a justification for
that value still has to come from the sampling design.
\citep{exec13_java_tetrad}

The missing-data policy also changes what is tested. The constructor can
fail on missing data, delete rows listwise, use an EM covariance route,
or use a testwise route. Its documented legacy behavior uses testwise
deletion with a warning when no explicit policy is supplied. Different
tests can then use different observed subsets. For example, a funding
variable available only after a customer responds makes the complete
cases a selected population. The proposed adapter must report the
selection and policy; reporting only the nominal input row count would
conceal a change in the population supporting the graph. Multiple
imputation is explicitly delegated to a wrapper rather than implemented
by a single Fisher-Z test. \citep{exec13_java_tetrad}

Failure semantics deserve equal attention. The inspected
\texttt{checkIndependence} method catches a singular-matrix exception
and returns dependence with \(p=0\). A different public path,
\texttt{getResult}, propagates the singular-matrix exception. The main
calculation also rejects inadequate degrees of freedom in its ordinary
finite-correlation branch. Consequently an apparent strongest possible
dependence decision can arise from a numerical failure path. A banking
example is the inclusion of both total balances and all their exact
accounting components: the singularity says something about the
representation, not necessarily about evidence for a causal edge.
\citep{exec13_java_tetrad}

The appropriate integration design is a typed test result containing the
statistic, conditioning variables, effective observations and a
computation status. If the chosen package method does not expose whether
a zero value came from a singularity, the adapter must select an API
that does, instrument a qualified wrapper, or decline that stronger
interpretation. Checking covariance rank once at input is useful but
does not reproduce every numerical condition inside every conditional
test. A specification that promises complete failure provenance from the
unmodified method would exceed the inspected interface.

\hypertarget{why-conditioning-on-every-neighbour-is-not-a-safe-shortcut}{%
\subsection{Why conditioning on every neighbour is not a safe
shortcut}\label{why-conditioning-on-every-neighbour-is-not-a-safe-shortcut}}

Aitiotita exposes the correlation measure, arrow-finding component,
graph constraints, maximum conditioning size, and stable, conservative
and acyclicity settings separately. Its inspected PC skeleton routine
adds a consequential shortcut once the conditioning depth exceeds two.
For each remaining pair it tests conditioning on the union of their
neighbours. If dependence persists, it marks the link undeletable before
enumerating the smaller subsets at that depth. This rule would need the
implication that failure of the largest conditioning set rules out every
smaller separator. Conditional independence does not have that
monotonicity property. \citep{exec13_java_aitiotita}

An exact project example makes the issue visible in a banking setting.
Let \(Z_1,Z_2,Z_3\) be three independent baseline funding factors, \(X\)
a rate concession, \(Y\) a balance response, and \(C\) a later
reconciliation score. For this diagnostic, deliberately stipulate no
causal effect of \(X\) on \(Y\). Use independent, centered,
unit-variance Gaussian inputs and \[
 X=Z_1+Z_2+Z_3+\varepsilon_X,\qquad
 Y=Z_1+Z_2+Z_3+\varepsilon_Y,\qquad
 C=X+Y+\varepsilon_C.
\] The shared baselines confound the observational relation, and the
later score is a collider. Conditioning on all three baselines leaves
independent residuals \(\varepsilon_X,\varepsilon_Y\), so \[
 \operatorname{cov}(X,Y\mid Z_1,Z_2,Z_3)=0.
\] After those baselines are held fixed, conditioning additionally on
\(C\) conditions on \(\varepsilon_X+\varepsilon_Y+\varepsilon_C\).
Gaussian conditioning gives \[
 \operatorname{cov}(X,Y\mid Z_1,Z_2,Z_3,C)
   =0-\frac{1\cdot1}{3}=-\frac13,
 \qquad
 \operatorname{var}(X\mid Z_1,Z_2,Z_3,C)=\frac23.
\] The variance for \(Y\) is also \(2/3\), hence the partial correlation
is \(-1/2\). Adding information reopens a path that the baseline
separator had closed. \citep{exec13_java_examples}

This example is deliberately constructed to reach depth three. If \(k\)
baseline factors remain unconditioned, then the residual covariance of
\(X,Y\) is \(k\), each residual variance is \(k+1\), their residual
covariance with \(C\) is \(2k+1\), and the residual variance of \(C\) is
\(4k+3\). Conditioning on \(C\) therefore changes that covariance to \[
 k-\frac{(2k+1)^2}{4k+3}=-\frac{k+1}{4k+3}.
\] Every candidate set of size at most two, drawn from
\(\{Z_1,Z_2,Z_3,C\}\), leaves nonzero covariance. The exact calculation
enumerates all eleven such subsets; the three baseline variables are the
first separator. The all-neighbour set includes \(C\) and produces
dependence instead. The population calculation thus invalidates the
general shortcut even before finite-sample testing errors are
considered. It does not claim a compiled Aitiotita run or identify the
behavior of every delegated component.
\citep{exec13_java_examples,exec13_java_aitiotita}

For adoption, the useful next experiment is narrowly specified. Supply
an exact independence oracle for this six-variable law, verify the
retained graph and separating sets, and compare a version with the
shortcut removed. Inspect the graph mutation, correlation and
orientation implementations as part of that experiment. If the shortcut
is necessary for a specialized admissible model class, state and prove
the required restriction. Until then, Aitiotita is a source for
comparison and possible component reuse, with direct adoption of this
path held pending that review. Its Apache 2.0 license does not establish
the statistical correctness of the path. \citep{exec13_java_aitiotita}

\hypertarget{a-classifier-graph-and-a-funding-intervention}{%
\section{A classifier graph and a funding
intervention}\label{a-classifier-graph-and-a-funding-intervention}}

\hypertarget{wekas-useful-role-and-its-causal-boundary}{%
\subsection{Weka's useful role and its causal
boundary}\label{wekas-useful-role-and-its-causal-boundary}}

Weka's Bayesian-network classifier provides a useful alternative
precisely because its objective is explicit in the source. The inspected
\texttt{BayesNet} first removes records with missing class, normalizes
the data, builds a structure, and estimates conditional probability
tables. The normalization imports a supervised discretizer and a
replacement filter for missing values. The default search object is K2.
The common search class initializes the maximum parent count to one and
the naive-Bayes option to true; that initialization adds arrows from the
class variable to every other attribute. \citep{exec13_java_weka}

Suppose the class is twelve-month deposit retention and the attributes
include income recorded before an offer. An arrow from retention to
prior income can be an entirely legitimate predictive factorization. It
is not evidence that intervening on retention changes previously
recorded income. Discretizing continuous attributes with information
from the class similarly serves the classifier objective; it requires a
separate argument before becoming an unsupervised causal-discovery
transformation. Rows removed because the class is missing also change
the analyzed population. The proposed system should retain Weka as a
predictive comparator with an explicit preparation pipeline, and reject
automatic conversion of its default graph into a causal diagram.
\citep{exec13_java_weka}

Selecting an algorithm whose name includes conditional independence does
not remove these questions. In the inspected \texttt{CISearchAlgorithm},
the test compares the local score of \(X\) with parents \(Z\) to the
score with \(Y\) added as a parent. It declares independence if the
extra parent fails to improve the score. The reviewed local-score class
defaults to a Bayesian score, and the IC search's maximum conditioning
cardinality is initialized to two. These source observations do not
describe a Fisher-Z test with an interchangeable significance level. An
optional Markov-blanket classifier correction is another reason to
retain the complete search configuration. \citep{exec13_java_weka}

The inspected test class inherits a general classifier test and returns
the default \texttt{BayesNet}. This is useful evidence about the testing
interface, but it is not a campaign measuring correct causal
orientation, identification or effect coverage. The current official
Weka source must also be acquired before extrapolating the old mirror's
defaults to a current proposal. The mirror's GPL text and Java 8 build
profile are recorded facts about those bytes; neither proves successful
installation in the proposed Java runtime. \citep{exec13_java_weka}

\hypertarget{a-matrix-transpose-can-survive-an-acyclicity-check}{%
\subsection{A matrix transpose can survive an acyclicity
check}\label{a-matrix-transpose-can-survive-an-acyclicity-check}}

Consider a second project diagnostic with liquidity pressure \(L\), rate
concession \(R\), retained balances \(B\), and net funding cost \(K\).
Use centered, independent disturbances and normalized model units: \[
 R=0.5L+\varepsilon_R,\qquad
 B=0.3L+0.8R+\varepsilon_B,\qquad
 K=0.1R-0.2B+\varepsilon_K.
\] The coefficients are stipulated examples. The direct rate cost is
offset by the reduction in alternative funding cost associated with
retained balances. An intervention increasing \(R\) by one model unit
changes balances by \(0.8\) and changes cost by \(0.1-0.2(0.8)=-0.06\).
This is a total intervention response under the specified equations, not
an observational regression result. \citep{exec13_java_examples}

With rows denoting sources and columns targets, in the order
\((L,R,B,K)\), the coefficient matrix is \[
 W=\begin{pmatrix}
 0&0.5&0.3&0\\
 0&0&0.8&0.1\\
 0&0&0&-0.2\\
 0&0&0&0
 \end{pmatrix},\qquad x=W^{\mathsf T}x+\varepsilon.
\] For a surgical intervention, replace the treatment equation, remove
its incoming coefficients and solve the remaining linear system.
Treating \(W^{\mathsf T}\) as though it used the same source-row
convention reverses all edges. The new graph is still acyclic and has
the same number of links, but \(R\) no longer has a directed path to
\(B\) or \(K\). Its modeled cost response is zero. An adapter test that
checks only shape, diagonal zeros, edge counts and cycles would accept
this scientifically wrong conversion. \citep{exec13_java_examples}

There is a second independent failure. If absolute weights are treated
as coefficients, the cost response becomes \(0.1+0.2(0.8)=0.26\),
reversing the decision implied by the original example. Tetrad's
inspected DAGMA conversion uses absolute weights to decide which
unweighted edges to retain; that is not a promise to return a signed
linear model. The adapter must preserve any available coefficient matrix
separately and must not reconstruct coefficients from the graph.
Tetrad's directed-edge constructor places a tail at the source and an
arrowhead at the target, while its endpoint type also distinguishes
circles. Endpoint semantics and numerical coefficients both require
explicit interchange contracts.
\citep{exec13_java_examples,exec13_java_tetrad}

Units add another test. If \(x'_i=d_i x_i\), where each \(d_i>0\) is a
unit conversion, then the source-row coefficient becomes \[
 W'_{ij}=\frac{d_j}{d_i}W_{ij}.
\] An absolute threshold applied to those coefficients can retain a
different graph after expressing deposits in dollars instead of
thousands of dollars. Standardization can fix a declared computational
convention, but the adapter must retain its fitted transformation and
translate intervention responses back to meaningful units. The proposed
conformance example should therefore include an asymmetric signed
matrix, a variable permutation, a unit change, and explicit DAG versus
CPDAG output. These are separate tests of direction, identity, scale and
retained ambiguity, not repeated checks of one serializer.

\hypertarget{functional-parity-restricted-derivatives-and-adoption-choices}{%
\section{Functional parity, restricted derivatives and adoption
choices}\label{functional-parity-restricted-derivatives-and-adoption-choices}}

\hypertarget{an-implementation-name-does-not-establish-numerical-parity}{%
\subsection{An implementation name does not establish numerical
parity}\label{an-implementation-name-does-not-establish-numerical-parity}}

The inspected DirectLingam helper supplies a small but revealing parity
question. After standardizing its input it calculates \[
 b_1=\frac1n\sum_i\frac{x_i^2}{2},\qquad
 b_2=\frac1n\sum_i x_i\exp(-x_i^2/2),
\] and uses these moments in an entropy-style expression. For unit
population variance, \(b_1=1/2\) regardless of the distribution's shape.
For a sample containing symmetric positive-negative pairs, the odd
second ingredient is zero. With sample-variance standardization the
common first ingredient becomes \((n-1)/(2n)\), so equal-length
symmetric samples still share it. This is an observation about the exact
source expression, which is more specific than its accompanying
description of a log-cosh-style approximation.
\citep{exec13_java_tetrad}

For example, standardizing \((-3,-1,1,3)\) and \((-1,-1,1,1)\) gives
fourth moments \(41/25\) and \(1\), respectively, under population
normalization. The two ingredients above are nevertheless identical. The
reproduced project calculation establishes that this helper cannot
distinguish these shapes from those ingredients alone. It does not
establish the selected order for a full dataset, the intended reference
expression, or a failure frequency of the complete algorithm. Those
stronger questions require authentic reference derivations and matched
execution, including residualization and parent selection. The proposed
adoption decision is to investigate this path before claiming parity
with another implementation bearing the same algorithm name.
\citep{exec13_java_examples,exec13_java_tetrad}

\hypertarget{banjo-preserve-the-original-rights-and-the-application-boundary}{%
\subsection{Banjo: preserve the original rights and the application
boundary}\label{banjo-preserve-the-original-rights-and-the-application-boundary}}

The inspected Banjo lead is a derivative whose README describes dynamic
gene-network analysis. It constructs pseudotime from single-cell
measurements, uses lag one, a stated discretization policy and a maximum
parent count, and selects gene subsets using supplied gold-standard
connectivity. The selected greedy-search source includes time, iteration
and restart termination. These are descriptions of an application and
its implementation choices; they are not banking benchmarks or evidence
that an inferred order is the recorded time of a bank action. Selecting
variables with reference to supplied graph truth would also have to be
disclosed when interpreting that application's evaluation.
\citep{exec13_java_banjo}

The source notice preserves Duke University and Alexander Hartemink's
rights. It says that non-commercial use may be licensed under a
non-commercial license and directs commercial users to contact the
rights holders. Public availability of the derivative therefore does not
support classifying it as an unrestricted open-source component. The
original official distribution, full terms and derivative changes remain
necessary for a reuse decision. Until those are established, the
proposal should compare the documented search family while holding
redistribution and commercial adoption. Rewriting only an interface
would not resolve the rights of copied implementation code.
\citep{exec13_java_banjo}

\hypertarget{what-to-integrate-what-to-retain-as-a-baseline-and-what-to-investigate}{%
\subsection{What to integrate, what to retain as a baseline, and what to
investigate}\label{what-to-integrate-what-to-retain-as-a-baseline-and-what-to-investigate}}

Tetrad is the leading candidate in this inspected Java subset for a
narrow PC/FGES integration because the chosen paths expose graph
semantics, tests, scores and configuration separately. The proposed
first implementation should use a pinned Java process with an explicit
request and response schema. This keeps its runtime and dependencies
identifiable and avoids immediately maintaining a second numerical port.
It does not remove GPL obligations or settle which runtime is required:
the inspected root compiler setting and bridge README name different
Java requirements. The actual dependency and bridge combination needs
its own build check. \citep{exec13_java_tetrad}

Weka should be retained for a separately labeled predictive comparison,
with current official sources obtained before promising current-release
behavior. Aitiotita offers a permissively licensed alternative whose
separation of components may be useful, but the conditioning shortcut
and delegated paths need focused examination before direct adoption. The
Banjo derivative supplies an additional search-family lead, with
original rights and technical identity still unresolved. A source port
is justified only when a specific integration constraint outweighs the
cost of preserving and validating the mathematical behavior. It is not
the default response to having several implementation languages.
\citep{exec13_java_weka,exec13_java_aitiotita,exec13_java_banjo}

The next prototype comparison should separate deterministic conformance
from statistical assessment. Deterministic cases check the funding
matrix, endpoint roundtrips, oracle separating sets, duplicate
accounting columns and typed numerical failures. Statistical cases
subsequently vary sample size, minimum signal, graph density,
non-Gaussianity, missingness and departures from the declared graph
class, with common generated truth and fixed replication budgets. A
single seeded package test can support a regression check; it cannot
estimate whole-graph error rates. The inspected Tetrad tests include
oracle-graph fixtures and a seeded chain/collider simulation, and one
test expressly calls its expected graph a regression snapshot. They
motivate these distinct levels of validation rather than replacing them.
\citep{exec13_java_tetrad}

Scale should be assessed at the same scientific boundary. A skeleton
search over \(p\) variables with at most \(d\) neighbours and depth
bound \(k\) can face a workload proportional to
\(p^2\sum_{j=0}^{k}\binom{d-1}{j}\) candidate tests before accounting
for pruning and the cost of each test. That expression is a planning
envelope, not a runtime measurement or a tight bound for every
algorithm. Dense graphs and expensive conditional tests can dominate the
result. Thread settings must be recorded from the selected code path:
the inspected FGES field and constructor initialize one worker even
though a nearby comment says all processors, and a setter can change
that value. Observed runtime, peak memory, interrupted requests and
retained failures should be measured only after a specific runtime and
workload are chosen. \citep{exec13_java_tetrad}

The resulting proposal is actionable without treating source review as a
finished product. It identifies a first Java boundary, preserves Weka's
predictive role, gives a concrete mathematical challenge for Aitiotita,
states the Banjo rights dependency, and defines the tests that would
distinguish a successful adapter from an attractive but scientifically
altered output. \label{exec13:java-discovery-end}


\section{Python discovery: choosing a computation that preserves the question}
\label{exec13:python-discovery}

A funding analyst can obtain three apparently plausible answers from the same
table: a partially directed graph, a matrix of regression coefficients, and a
distribution over neural structural models. Their common Python interface does
not make these answers interchangeable. The first can preserve directions that
the data have not resolved. The second can describe one selected ordering. The
third can express uncertainty conditional on a generative model, its prior and
an approximation to its posterior. A product that converts all three into a
single directed network would discard precisely the distinctions needed to
decide whether changing a quoted rate is supported by evidence.

The comparison below examines nine implementation families at identified source
revisions preceding the 7 October 2026 cutoff: causal-learn, pgmpy, LiNGAM,
NOTEARS, DAGMA, gCastle, CausalNex, Causica and the CausalDiscoveryToolbox (CDT).
It develops selected static discovery paths, rather than attributing every
algorithm advertised by a package to the code inspected. These are proposed
reuse choices supported by primary implementation and documentation. The
financial and mathematical examples are reproducible project calculations;
they are not fitted outputs from these engines. Broader Python search coverage,
some delegated implementations and empirical comparisons remain open.

\subsection{The funding calculation establishes the matrix contract}

Consider three variables measured over a fixed decision horizon. Let \(F\) be a
change in a quoted funding rate, in percentage points, \(R\) the change in
retained funding, in millions of US dollars, and \(C\) the corresponding change
in contribution margin, in the same monetary units. A hypothetical structural
model is
\begin{equation}
 F=\varepsilon_F,\qquad R=4F+\varepsilon_R,\qquad
 C=-F+\tfrac12R+\varepsilon_C.
 \label{exec13:pydisc-funding}
\end{equation}
The disturbances are centered and mutually independent, and intervention on
\(F\) replaces its first equation. These assumptions define the illustration;
the coefficients are not estimated bank responses. The direct margin cost of
raising the rate is offset by the margin earned on retained funding. A rate
increase of \(0.25\) percentage points therefore changes expected retention by
USD~1 million and expected margin by USD~0.25 million.

For a data matrix whose columns have order \((F,R,C)\), write the row convention
as \(X=XW+E\), with \(W_{ij}\) the coefficient on the edge \(i\to j\). The
corresponding column-vector structural convention is \(x=Bx+e\), where
\(B=W^\top\). In this example,
\begin{equation}
 W=\begin{pmatrix}0&4&-1\\0&0&1/2\\0&0&0\end{pmatrix},\qquad
 T=I+W+W^2=(I-W)^{-1}
  =\begin{pmatrix}1&4&1\\0&1&1/2\\0&0&1\end{pmatrix}.
\end{equation}
Thus \(T_{FC}=1=-1+4(1/2)\) is the total margin response per percentage point.
The matrix identity holds because \(W^3=0\). More generally, a finite DAG with
\(d\) nodes has a nilpotent weighted adjacency matrix and a finite path sum
\(I+W+\cdots+W^{d-1}\). This algebra concerns a supplied linear structural
model; it is not a discovery or identification theorem.

The reference NOTEARS and DAGMA linear losses predict columns through \(XW\).
The inspected LiNGAM adjacency routine instead assigns coefficients to
\texttt{B[target,predictors]} and its path helper multiplies entries in that
convention \citep{exec13_pydisc_notears,exec13_pydisc_dagma,exec13_pydisc_lingam}.
An adapter that interprets LiNGAM's matrix as \(W\) reports zero for the
\(F\to C\) response in this example and invents the reversed path response.
Both matrices are acyclic, both have identical dimensions, and both pass a
check that merely counts edges. Node order, direction convention and units
must therefore be carried as data, verified before any intervention query.

Endpoint matrices require a separate decoder. In causal-learn's PC/GES return
contract, entries \((G_{ij},G_{ji})=(-1,1)\) denote \(i\to j\); a pair of
minus ones denotes an undirected edge. These entries are endpoint marks, not
negative and positive causal coefficients \citep{exec13_pydisc_causallearn}.
Treating them as weights would invent both magnitudes and interventions. The
proposed common result must distinguish a weighted DAG, a partially directed
graph and a posterior over structural models before it exposes any common
plotting or effect-evaluation interface.

\subsection{A finite landscape of complementary mechanisms}

Table~\ref{tab:exec13-pydisc} records a proposed role for each inspected family.
The recommendations depend on the scientific target, not on a single numerical
ranking. A package can be useful as a comparison engine without being the
default dependency, and a method can be worth retaining even when it cannot
resolve the bank's current ambiguity.

\begingroup\small
\begin{longtable}{p{0.16\textwidth}p{0.36\textwidth}p{0.40\textwidth}}
\caption{Selected Python discovery paths and proposed reuse. These are source
comparisons, not measured engine rankings.}\label{tab:exec13-pydisc}\\
\toprule Family & Inspected mechanism & Proposed role and condition \\
\midrule\endfirsthead
\toprule Family & Inspected mechanism & Proposed role and condition \\
\midrule\endhead
causal-learn & PC endpoint graph; GES with selectable scores & Integrate selected
paths with explicit CI, score, collider and graph-type contracts. \\
pgmpy & PC; equivalence search with forward, backward and turning phases &
Integrate the pinned graph API; request partially directed output explicitly. \\
LiNGAM & DirectLiNGAM order selection; coefficient and total-path helpers &
Integrate within a declared linear, non-Gaussian, causally sufficient model;
preserve the target-row convention. \\
NOTEARS & Linear loss, exponential acyclicity and sparsity penalty & Retain
the compact reference as a numerical baseline before considering a port. \\
DAGMA & Linear loss with log-determinant acyclicity and domain control & Retain
a separate baseline; check domain, final topology, input mutation and constraints. \\
gCastle & Selected PC and polynomial-constraint NOTEARS paths & Retain a
matched alternative with method-specific semantics and thresholds. \\
CausalNex & StructureModel and NOTEARS variants with tabu constraints & Defer
default adoption; evaluate an isolated environment for a concrete consumer. \\
Causica & Learned structural functions, noise and graph distribution & Retain
an experimental generative-model option; evaluate effects conditionally. \\
CDT & Selected PC/GES wrappers invoking R packages & Defer default adoption
and porting until delegated versions, ordering and terms are resolved. \\
\bottomrule
\end{longtable}
\endgroup

The source evidence for these rows is the selected PC/GES implementations
\citep{exec13_pydisc_causallearn,exec13_pydisc_pgmpy}, the functional and
continuous implementations
\citep{exec13_pydisc_lingam,exec13_pydisc_notears,exec13_pydisc_dagma,exec13_pydisc_gcastle},
and the framework and wrapper paths
\citep{exec13_pydisc_causalnex,exec13_pydisc_causica,exec13_pydisc_cdt}.
\texttt{castle} is the import namespace of gCastle; it is not an additional
independent implementation. The repository \texttt{cdt15/lingam} and CDT are
distinct projects despite their similar names. Repository names and import
aliases should be resolved before counting competitors or estimating the
cost of integration.

\subsection{Constraint and score methods preserve different qualifications}

The selected causal-learn PC interface exposes the significance threshold,
conditional-independence test, stable skeleton option, collider rule, conflict
priority, background knowledge and conditioning-set bound. It also has a
separate missing-value route. These arguments control substantive choices in
the graph construction. They should appear in the result's explanation whenever
they affect the decision, rather than disappearing into an undifferentiated
``PC'' label \citep{exec13_pydisc_causallearn}.

For its ordinary Fisher test, the source computes a partial correlation from
an inverse correlation submatrix, then forms
\begin{equation}
 Z=\sqrt{n-|S|-3}\,\bigl|\operatorname{atanh}(\widehat r_{ij\mid S})\bigr|,
 \qquad p=2\{1-\Phi(Z)\}.
\end{equation}
Here \(S\) is the conditioning set, \(n\) the row count and \(\Phi\) the
standard normal distribution function. The Gaussian independent-row sampling
model motivates this reference approximation; correlation zero does not imply
conditional independence for an arbitrary non-Gaussian distribution. The code
raises an error for a singular correlation submatrix and clips a computed
correlation at the numerical boundary. Neither clipping nor returning a
number establishes adequate information. A sensible adapter must reject
nonpositive \(n-|S|-3\), constant columns and singular designs with an
explanation, not translate failure into an independence finding.

Repeated observations from one account also challenge this reference law. If a
correlation of \(0.1\) is supplied with 100 independent units, the displayed
statistic is approximately \(0.988\). Substituting 1,000 duplicated rows gives
approximately \(3.168\), although the duplicated data contain no new
independent information. The example compares the formula under two declared
row counts; it does not propose a universal effective-sample-size correction
\citep{exec13_pydisc_examples}. A panel or time-series request must use a
sampling argument appropriate to that dependence, or be declined by this
static interface.

The inspected pgmpy PC path defaults to a partially directed graph and limits
conditioning sets to five variables unless configured otherwise. Its optional
DAG is one completion of the learned partial graph. Its documented background
knowledge includes required and forbidden directions, while collider conflicts
can still generate a warning. Consequently the proposed adapter should check
the returned graph against declared constraints instead of treating the mere
presence of an expert-knowledge argument as proof that every input condition
was satisfied \citep{exec13_pydisc_pgmpy}. A five-variable search limit is a
computational scope choice, not an assertion that larger separators cannot
matter in the bank's model.

Score methods make a different tradeoff. causal-learn GES accepts several
local scores, including Gaussian BIC, discrete BDeu and regression-based
alternatives. Its inspected BIC routine computes residual variance and can use
a pseudoinverse or a small positive replacement for a nonpositive numerical
variance. pgmpy's inspected GES performs forward, backward and turning moves,
uses a minimum improvement threshold, and converts its final equivalence graph
to a DAG only if requested
\citep{exec13_pydisc_causallearn,exec13_pydisc_pgmpy}. A larger score in one
engine cannot be compared directly with a number from another engine without
matching sign, constants, likelihood, penalties and preprocessing. The
proposal should compare graph/effect consequences under matched assumptions,
not select an engine because its default score is numerically larger.

\subsection{Functional restrictions can orient graphs, but must be stated}

LiNGAM is attractive when the banking question justifies a linear acyclic model
with independent non-Gaussian disturbances and no unobserved confounders. Those
are explicit premises in the original work, including independent repeated
observations from the same model \citep{exec13_pydisc_shimizu2006}. The original
paper describes an ICA route; the current DirectLiNGAM implementation inspected
here selects an ordering through residual comparisons, then estimates
coefficients using that order. The original result should not be quoted as
if it had independently verified every modern option and integration.

The selected helper standardizes variables for an adaptive-Lasso pruning step,
uses a BIC criterion, and refits the retained predictors in original units.
The separate total-effect helper sums directed path products and checks that
variables along those paths satisfy its continuous-variable condition
\citep{exec13_pydisc_lingam}. These are valuable implementation details for a
funding adapter: they explain where a sparse edge set comes from and how the
reported monetary slope is recovered. They do not show that bank disturbances
are independent or non-Gaussian. Heavy tails in an observed balance series
alone do not establish the required disturbance model after conditioning on
the structural parents.

The observational ambiguity introduced earlier provides an exact adverse case.
For a standardized pair, consider two centered Gaussian models,
\begin{align}
 \mathcal M_{\to}:&\quad X=U,\quad Y=0.6X+V,
     &\operatorname{var}(U)&=1,&\operatorname{var}(V)&=0.64,\\
 \mathcal M_{\leftarrow}:&\quad Y=\widetilde U,\quad X=0.6Y+\widetilde V,
     &\operatorname{var}(\widetilde U)&=1,&\operatorname{var}(\widetilde V)&=0.64.
\end{align}
Disturbances are independent within each model. Both laws have covariance
\(\left(\begin{smallmatrix}1&0.6\\0.6&1\end{smallmatrix}\right)\), so the
complete observational distributions coincide, not just a selected correlation.
The effect of intervening on \(X\) is nevertheless \(0.6\) in the first model
and zero in the second. Within this declared two-model class the effect set is
\(\{0,0.6\}\); the interval \([0,0.6]\) is its convex hull, not an assertion
that every intermediate effect belongs to the two-model class.

No algorithm using observations alone can distinguish these two Gaussian laws
reliably in both worlds. A non-Gaussian identifying premise, a justified
temporal direction, or an intervention changes the available information.
An optimizer's ability to return one orientation does not. The proposed
DirectLiNGAM integration therefore retains this example as a refusal or
assumption-failure case, rather than interpreting a selected order as empirical
proof of its own premises.

\subsection{Continuous acyclicity is a numerical condition with a domain}

Continuous methods replace a combinatorial acyclicity constraint by a smooth
one. Let \(S=W\circ W\), the entrywise square of \(W\). NOTEARS uses
\begin{equation}
 h_{\exp}(W)=\operatorname{tr}(\exp S)-d
             =\sum_{k\geq1}\frac{\operatorname{tr}(S^k)}{k!}.
\end{equation}
Because \(S\) is nonnegative, its diagonal powers accumulate nonnegative
closed-walk products. A directed cycle creates a positive term; in a DAG all
such terms vanish. This is the exact-zero characterization developed in
\citet{notears2018}. Its numerical implementation uses tolerances, iteration
limits and thresholding, each of which must be distinguished from the exact
statement.

The source inspected for gCastle's class named \texttt{Notears} instead uses
the polynomial expression
\begin{equation}
 h_{\rm poly}(W)=\operatorname{tr}\{(I+S/d)^d\}-d.
\end{equation}
The reference NOTEARS linear implementation uses a matrix exponential. The
two expressions share an exact-zero acyclicity property but have different
values, derivatives and numerical scaling away from zero. gCastle also keeps
the weighted estimate separately from a binary matrix formed with the strict
comparison \(|W_{ij}|>\tau\), whereas the reference removes weights satisfying
\(|W_{ij}|<\tau\). At equality, the reference retains the edge and gCastle's
binary output does not \citep{exec13_pydisc_gcastle,exec13_pydisc_notears}.
A benchmark that silently compares one weighted matrix with another binary
matrix can mistake a representation choice for a scientific result.

DAGMA uses the log-determinant characterization
\begin{equation}
 h_{\log,s}(W)=-\log\det(sI-S)+d\log s,
 \qquad s>\rho(S),
\end{equation}
where \(\rho(S)\) is the spectral radius. The domain ensures that \(sI-S\) is
a nonsingular M-matrix. The domain condition is part of the theorem, not an
optional numerical detail \citep{exec13_pydisc_dagma2022}. For example, within
this domain one can expand the log determinant as
\(\sum_{k\geq1}\operatorname{tr}(S^k)/(k s^k)\); its nonnegative cycle terms
again explain the zero condition. Computing a log absolute determinant outside
the domain does not justify the same conclusion.

The inspected linear DAGMA implementation uses inverse-matrix checks and
backtracking in its optimization. Its documentation explicitly allows the
possibility that the finite output is cyclic and states that requested included
edges are not guaranteed to survive. It also centers the supplied array in
place and saves a final residual before thresholding
\citep{exec13_pydisc_dagma}. These findings motivate copying the input,
retaining the full optimization configuration, and recomputing topology and
residual diagnostics on the actual returned graph. They establish a proposed
adapter requirement; they do not establish a measured failure rate.

A small exact counterexample illustrates why a tolerance cannot replace a
topological check. Set
\(W=\left(\begin{smallmatrix}0&\epsilon\\\epsilon&0\end{smallmatrix}\right)\)
with \(\epsilon=0.005\). This is a two-cycle. Nevertheless,
\begin{align}
 h_{\exp}(W)&=2\{\cosh(\epsilon^2)-1\}\simeq6.25\times10^{-10},\\
 h_{\rm poly}(W)&=\epsilon^4/2=3.125\times10^{-10},\\
 h_{\log,1}(W)&=-\log(1-\epsilon^4)\simeq6.25\times10^{-10}.
\end{align}
All are below \(10^{-8}\). A threshold of \(0.001\) retains both edges. The
reference default threshold of \(0.3\) removes both; this example is therefore
not a claim that a default solver necessarily returns this cycle. It shows
why the numerical residual, thresholded support and scientific effect are
three distinct objects \citep{exec13_pydisc_examples}.

\subsection{Units and centering can change the selected orientation}

In the preceding two-model Gaussian example, the population least-squares
loss \(\tfrac12\sum_j\operatorname{var}(\varepsilon_j)\) is \(0.82\) in
both orientations, and the coefficient \(\ell_1\) penalty is \(0.6\) in both.
Both graphs are acyclic. Thus neither perfect acyclicity nor this matched
population objective distinguishes the two causal effects. Scores that add
other model restrictions may distinguish them, but those restrictions must be
declared and defended.

Now express \(Y\) in units ten times as large numerically, \(Y'=10Y\), while
leaving \(X\) unchanged. The forward coefficient becomes \(6\), the reverse
coefficient becomes \(0.06\), and their residual-variance sums become \(65\)
and \(100.64\), respectively. The corresponding least-squares losses are
\(32.5\) and \(50.32\). The physical observational ambiguity is unchanged,
but the unweighted loss and sparsity tradeoff have changed
\citep{exec13_pydisc_examples}. This matters when a funding dataset mixes
percentage points, currency, duration and account counts. Standardization is a
modeling decision that must be recorded and reversed when reporting monetary
responses; it does not independently supply causal identification.

CausalNex's selected unpenalized NOTEARS path applies its squared-error
criterion to the supplied numeric array without the centering performed in the
reference NOTEARS least-squares path. It provides tabu edges and prohibited
parent/child nodes as optimization bounds, and its StructureModel records edge
origins. The same graph object has an explicit operation that repeatedly drops
the smallest absolute-weight edge until a DAG remains, with a warning about
different units \citep{exec13_pydisc_causalnex}. Dropping an edge changes the
model used by an intervention query. It must not be treated as invisible file
cleanup, and the thresholded model should be assessed on its own assumptions.

The pinned CausalNex package declares Python versions from 3.8 to below 3.11.
That declaration does not prove installation failure in every newer
environment, but it is sufficient to reject an unqualified promise that the
current common Python environment can adopt it directly. An isolated process
is a possible design if a consumer values its StructureModel workflow enough
to justify the additional environment. For the simpler linear reference
comparison, wrapping NOTEARS directly is a substantive alternative that avoids
introducing that framework dependency merely to obtain a DAG matrix.

\subsection{Causica connects discovery and effects through a learned model}

Causica's inspected DECI path jointly learns a distribution over graphs,
functional relationships and noise. Its training step samples a relaxed graph,
evaluates a data log density, incorporates the graph distribution's entropy
and a graph prior, and adds an acyclicity penalty. Suppressing implementation
scaling details, the statistical part has the form
\begin{equation}
 -\mathbb E_{q_\phi(G)}\log p_\theta(X\mid G)
 +\operatorname{KL}\{q_\phi(G)\,\|\,p(G)\}.
\end{equation}
Here \(q_\phi\) is an approximating graph distribution, \(\theta\) parameterizes
the structural functions and noise, and \(p(G)\) contains graph preferences.
The actual module uses minibatches, relaxed samples and scaled penalties;
this expression explains their statistical roles, not an assertion of exact
posterior computation. The inspected prior includes sparsity and an optional
weighted discrepancy from a supplied expert graph
\citep{exec13_pydisc_causica}.

For an assumed structural model, simulation after replacing an intervention
equation can estimate a model-implied ATE. The Causica documentation describes
this workflow, while the inspected test methods separately evaluate graph
reconstruction, observational likelihood, interventional likelihood and effect
errors when appropriate reference data are supplied. These distinctions are
useful for the proposed product. A model may predict observed balances well
while choosing the wrong intervention direction in an observationally ambiguous
case; graph uncertainty under \(q_\phi\) is conditional on the model family,
prior and approximation. It is not automatically a confidence set with
frequentist coverage for a bank effect.

In particular, an expert graph supplied by a language model should retain its
origin and uncertainty. A strong prior favoring a rate-to-balance direction can
concentrate a fitted graph distribution without providing new assignment
evidence. A suitable prototype would compare the same task under an empty
prior, a justified restricted prior and a deliberately wrong prior, preserving
the resulting change in intervention conclusions. Such a comparison is a
proposed experiment, not a result of this source review.

The exact revision also matters for capability claims. Its README states that
Rhino and DDECI were removed from the refactored codebase and points to an older
release. Their historical association with Causica cannot be used to claim
that the inspected DECI path currently implements general temporal or
latent-confounded inference. A missingness or latent-variable claim needs the
corresponding implementation and assumptions, examined in its own right.

\subsection{A Python wrapper can carry another language's scientific contract}

CDT's selected PC and GES classes illustrate why language should be recorded
at both the public interface and the computational backend. PC checks for
\texttt{pcalg}, \texttt{kpcalg} and \texttt{RCIT}; GES checks for
\texttt{pcalg}. The wrappers construct files, dispatch R templates and read
results into Python graph objects \citep{exec13_pydisc_cdt}. A Python package
pin alone therefore does not freeze the statistical engine, the R environment
or its graph encoding. A direct R adapter is a serious alternative if it
preserves the same semantics with fewer conversion layers.

There are two concrete reasons to defer a default CDT recommendation. First,
its GES documentation labels one option as categorical, while the code maps
the corresponding \texttt{int} key to \texttt{GaussL0penIntScore}. The selected
Python wrapper does not establish a discrete likelihood merely because of
that prose label. The R template and delegated implementation need inspection
before assigning the option a capability. The wrapper also relabels graph
nodes using sorted column names in a selected path, making a permutation test
with deliberately unsorted columns essential before banking use.

Second, the repository's saved license text is MIT, while the saved setup
metadata names Apache~2.0. Both observations are retained. Neither is a reason
to invent a single reconciled distribution license or to assume delegated
R packages inherit the wrapper's terms. This is a local adoption question,
not a reason to stop comparisons of unrelated Python implementations. The
other selected primary terms are Apache~2.0 for NOTEARS, DAGMA and the inspected
gCastle root; CausalNex supplies Apache terms and separate trademark text;
Causica, LiNGAM and the pinned pgmpy license supply MIT text. causal-learn's
MIT text was acquired on the cutoff date from a mutable-main snapshot and
should be checked at the technical pin before distribution. These observations
guide dependency review; they do not constitute a complete transitive-license
opinion.

\subsection{Selection and missingness change what a discovery table means}

A funding table may contain only customers who were successfully contacted,
retained an account or completed a follow-up. Treating the remaining rows as a
random sample can create edges even before choosing a discovery algorithm.
For an exact generated illustration, let \(F\) and \(N\) be independent
standard normal variables and observe only rows satisfying \(F+N>0\). Define
\(U=(F+N)/\sqrt2\) and \(V=(F-N)/\sqrt2\). Before selection, \(U\) and \(V\)
are independent standard normals. After selection, \(U\) is half-normal with
variance \(1-2/\pi\), while \(V\) still has variance one. It follows that
\begin{equation}
 \operatorname{cov}(F,N\mid F+N>0)=-1/\pi,\qquad
 \operatorname{corr}(F,N\mid F+N>0)=-\frac1{\pi-1}.
\end{equation}
The selected correlation is about \(-0.467\), despite population independence
\citep{exec13_pydisc_examples}. A graph of the selected table can correctly
describe its associations while misleading a population intervention query.

This example gives practical meaning to the distinction between a missing-value
option and an argument for selection recovery. A package-specific route must
state which response mechanism it handles and what observations identify its
correction. Imputation, complete-case analysis and a generative missing-data
model make different commitments. The proposed request should carry the
sampling unit, observation window, selection indicators, and relevant
conditioning events. If these are unavailable, the product can still support
exploratory associations and alternative diagrams, but its intervention
conclusion must retain the unresolved selection assumption.

\subsection{An adapter contract and an economical prototype choice}

The proposed static-discovery boundary first validates the question and data:
units, node order, information timing, row dependence, missingness, allowed
interventions and model family. It then configures one pinned engine and
retains its native result. A second step decodes that result into a typed graph
or structural-model object. Only a subsequent identification step may turn
that object into a causal effect or effect set. For a partially directed
result, selecting one completion is an explicit additional assumption; for a
neural SEM, reporting simulation output is conditional on the selected learned
model and uncertainty representation.

Useful deterministic acceptance cases follow directly from the comparison.
The funding model tests a transpose and a deliberate permutation of columns.
The two Gaussian orientations test whether an interface preserves unresolved
effects. The weak two-cycle tests support-level acyclicity independently of a
floating-point residual. A coefficient exactly on the threshold tests the
inequality convention. A singular CI input tests refusal, while duplicated
rows test whether the declared sampling unit is preserved. These cases check
the proposed boundary; they are not substitutes for a prospectively designed
empirical evaluation of discovery accuracy or uncertainty.

The initial prototype should therefore favor a small, interpretable set of
mechanisms: one PC/GES path that preserves equivalence, one functional path
with explicit non-Gaussian assumptions, and separate NOTEARS/DAGMA numerical
references. gCastle can test the cost of a broader algorithm catalog.
CausalNex and CDT should enter only when a specific consumer justifies their
additional dependency or conversion burden. Causica supplies a distinct
generative-model experiment, whose contribution should be judged on matched
intervention questions and negative cases rather than an attractive graph
alone. Porting a numerical core becomes worthwhile only when measured
deployment constraints or performance needs exceed the cost of validating a
second implementation.

These choices make the proposal concrete without assigning a fictional bank
ROI. Their economic comparison must include environment maintenance, source
and assumption review, translation errors, failed requests and the cost of
resolving consequential graph ambiguity. The known routes are exact interface
tests, scoped primary-source completion and prospectively designed public or
generated comparisons. The unresolved scientific question is whether the
information available for a particular bank decision supports the required
intervention claim. Adding another optimizer does not answer that question.
\label{exec13:python-discovery-end}


\section{Symbolic identification and numerical evaluation}

aGrUM's inspected \texttt{CausalModel} documentation permits a causal DAG with
an optional observational Bayesian network. Structural operations, including
backdoor/frontdoor queries and construction of do-calculus expressions, can
be available when conditional probability tables are absent. Operations that
read those tables then raise an unavailable-operation exception. This is a
useful separation: a funding query can be identified relative to a supplied
graph before a distribution has been fitted. It does not follow that the
effect can already be reported numerically.\cite{v12_sw_agrum}

The selected \texttt{CausalImpact} implementation first attempts appropriate
backdoor and frontdoor paths, then invokes a do-calculus routine. The inspected
dispatch supports a comparison of these routes; it is not a verification of
the full ID or IDC algorithm. A particularly consequential detail is the
handling of failures. A not-yet-implemented method and a hedge exception are
caught in different branches, but each produces a formula with no expression
tree. Their explanation strings differ. An adapter that reads only the boolean
\texttt{isIdentified} would erase the distinction between an unavailable
implementation and a causal nonidentification result.\cite{v12_sw_agrum}

These outputs should instead be represented as separate states:
\begin{center}
\begin{tabular}{p{.29\textwidth}p{.60\textwidth}}
\toprule
State & Required scientific meaning\\
\midrule
Identified expression & A functional has been obtained under the declared graph and assumptions.\\
Numerical result unavailable & The expression may exist, but distributions, estimators or numerical prerequisites are absent.\\
Method unsupported & This implementation cannot perform the requested method or graph operation.\\
Nonidentified in the model & The supplied causal model and information do not determine the requested functional, with the relevant reason retained.\\
Computational failure & The attempted operation failed; the failure supplies no causal conclusion.\\
\bottomrule
\end{tabular}
\end{center}

pgmpy's guide also separates identification from estimation and uses annotated
graph roles. Its adjustment interface returns a graph or collection of graphs
and a success indicator. Source inspection narrows what the guide alone would
suggest: the \texttt{minimal\_variance} branch raises
\texttt{NotImplementedError}, despite the guide's description of variance-optimal
adjustment. The minimal branch rejects multiple exposures or outcomes. The
selected implementation contains a graph-type TODO, so the broad graph-type
docstring does not by itself qualify every graph class.\cite{v12_sw_pgmpy}

An unsuccessful adjustment search is also narrower than failure of general
causal identification. The package guide discusses a frontdoor route as a
distinct alternative. The proposed product should retain the method attempted
and its applicability rather than mapping every false adjustment indicator to
``causal effect impossible.'' For a bank analyst, that distinction determines
whether the next action is to consider another identifying functional, obtain
additional information, or correct an unsupported request.

\section{Effect candidates, hulls and confidence sets}

Tetrad's inspected IDA constructor requires a legal partially directed acyclic
graph and continuous data. Its total-effect routine produces regression-based
candidates across eligible parent configurations. Candidate-specific exceptions
are caught and logged without a typed failed-candidate object in the returned
list. Consequently, the downstream interface needs enough information to
distinguish a genuinely small candidate set from a larger intended calculation
whose members partly failed.\cite{v12_sw_tetrad}

Three numerical objects must remain distinct. A candidate multiset records
effects attached to graph alternatives. Its convex hull fills all values
between the smallest and largest candidate. A confidence set is a random set
with a stated coverage interpretation under a sampling law. Neither the
candidate list nor its hull becomes a confidence set merely because it has
endpoints.

The inspected IDA distance helper illustrates the distinction. For project
values \(S=\{0,0.8\}\) and hypothetical truth \(\theta=0.4\),
\[
 d(\theta,S)=\min_{s\in S}|\theta-s|=0.4,
 \qquad d(\theta,\operatorname{hull}(S))=0.
\]
The selected helper uses the minimum--maximum interval, so it computes the
second quantity. Its empty-list branch returns the supplied true effect, which
could even be negative. These are static observations about that helper,
accompanied by project arithmetic; no Tetrad execution is claimed. They imply
that evaluation and transport must define their set object before choosing a
distance metric. An empty set of available estimates requires an explicit
unavailable or failed state, not an arbitrary scalar score.\cite{v12_sw_tetrad}

The initial adoption preference is therefore to evaluate the existing Java
engine through a narrow process boundary, with endpoint roundtrips and typed
failure handling. Porting an isolated helper would omit the graph semantics
that make its output meaningful. For aGrUM, the existing C++ core and Python
wrapper deserve comparison before any reimplementation. Its saved build files
declare version 3.2.0 and C++20; these are environment requirements to test,
not evidence that the proposed binary or wrapper already works.

